% Hafta 2 -- alıştırma gövdesi (ders notu ve ayrı alıştırma kağıdı tarafından \input edilir)
\section{Çözümlü Alıştırmalar}
Her soruyu önce kendiniz deneyin, sonra çözümü okuyun.

\begin{alistirma}
\(p\) ve \(p+9\) sayılarının ikisi de asal ise \(p+(p+9)\) toplamı kaçtır?
\end{alistirma}
\begin{cozum}
\(p\) tek bir asal olsaydı \(p+9\) çift ve \(2\)'den büyük olurdu, yani asal olamazdı. O halde \(p\) çift asaldır: \(p=2\).\\
\(p+9=11\) asal \(\checkmark\). \quad Toplam: \(2+11=13\).
\end{cozum}

\begin{alistirma}
\(1\)'den \(24\)'e kadar (\(1\) ve \(24\) dahil) olan tam sayılardan kaç tanesi \(10\) ile aralarında asaldır?
\end{alistirma}
\begin{cozum}
\(10=2\cdot 5\): sayı ne \(2\)'ye ne \(5\)'e bölünmeli.\\
\(2\)'nin katları: \(12\) tane; \ \(5\)'in katları: \(5,10,15,20\) \(\to\) \(4\) tane; \ ikisinin de katı (\(10,20\)): \(2\) tane.\\
\(2\)'ye veya \(5\)'e bölünenler: \(12+4-2=14\). \quad Aralarında asal olanlar: \(24-14=10\).\\
(Liste: \(1,3,7,9,11,13,17,19,21,23\).)
\end{cozum}

\begin{alistirma}
\(120\) sayısının asal olmayan kaç pozitif böleni vardır?
\end{alistirma}
\begin{cozum}
\(120=2^3\cdot 3\cdot 5\) \(\Rightarrow\) pozitif bölen sayısı \((3+1)(1+1)(1+1)=16\).\\
Asal bölenler: \(2,3,5\) \(\to\) \(3\) tane. \quad Asal olmayanlar: \(16-3=13\) \ (\(1\) de asal değildir, bu sayıya dahildir).
\end{cozum}

\begin{alistirma}
\(90\) sayısının pozitif bölenlerinden kaç tanesi çift sayıdır?
\end{alistirma}
\begin{cozum}
\(90=2\cdot 3^2\cdot 5\) \(\Rightarrow\) bölen sayısı \(2\cdot 3\cdot 2=12\).\\
Tek bölenlerde \(2\) çarpanı bulunmaz: \(3^2\cdot 5\)'in bölen sayısı \(3\cdot 2=6\).\\
Çift bölenler: \(12-6=6\) \ (\(2,6,10,18,30,90\)).
\end{cozum}

\begin{alistirma}
\(108\cdot n\) bir tam küp ise \(n\) en küçük hangi pozitif tam sayıdır?
\end{alistirma}
\begin{cozum}
\(108=2^2\cdot 3^3\). Tam küpte bütün üsler \(3\)'ün katı olmalıdır.\\
\(2^2\to 2^3\) için bir \(2\) gerekir; \(3^3\) zaten uygundur. \quad \(n=2\), \ \(108\cdot 2=216=6^3\).
\end{cozum}

\begin{alistirma}
Bir doğal sayı \(15\) ile bölündüğünde bölüm \(7\) oluyor. Bu sayı en az ve en çok kaçtır?
\end{alistirma}
\begin{cozum}
\(A=15\cdot 7+K=105+K\), \quad \(0\le K\le 14\).\\
En az: \(K=0\Rightarrow A=105\). \quad En çok: \(K=14\Rightarrow A=119\).
\end{cozum}

\begin{alistirma}
\(A\)'nın \(7\) ile bölümünde bölüm \(B\), kalan \(4\); \ \(B\)'nin \(3\) ile bölümünde bölüm \(C\), kalan \(2\)'dir. \(A\)'nın \(21\) ile bölümünden kalan kaçtır?
\end{alistirma}
\begin{cozum}
\(A=7B+4\) \ ve \ \(B=3C+2\). \ Yerine yazalım:
\[A=7(3C+2)+4=21C+14+4=21C+18\]
\(18<21\) olduğundan kalan \(18\)'dir. \ (Kontrol: \(C=0\Rightarrow B=2,\ A=18\).)
\end{cozum}

\begin{alistirma}
Pozitif bir tam sayının \(9\) ile bölümünde bölüm, kalanın \(2\) katıdır. Bu koşulu sağlayan en büyük sayı kaçtır?
\end{alistirma}
\begin{cozum}
Kalan \(K\), bölüm \(2K\): \ \(A=9\cdot 2K+K=19K\).\\
\(0\le K\le 8\); sayı pozitif olduğundan \(K\ge 1\). En büyük \(A\) için \(K=8\):
\[A=19\cdot 8=152\qquad(\text{kontrol: } 152=9\cdot 16+8,\ 16=2\cdot 8)\]
\end{cozum}

\kendinizoncesi
\section{Alıştırmalar (Kendiniz Çözün)}

\begin{enumerate}[label=\textbf{\arabic*.},leftmargin=*,itemsep=6pt]
\item \(30\) ile \(60\) arasında kaç asal sayı vardır?
\item \(187\) ve \(191\) sayıları asal mıdır?
\item \(a\) ve \(b\) asal sayılardır. \(a+b=31\) olduğuna göre \(a\cdot b\) kaçtır?
\item \(252\) sayısını asal çarpanlarına ayırınız. Kaç pozitif böleni vardır?
\item \(400\) sayısının kaç tam böleni vardır?
\item \(2^a\cdot 5^3\) sayısının \(20\) pozitif böleni vardır. \(a\) kaçtır?
\item \(18\) sayısının pozitif bölenlerinin toplamı kaçtır?
\item \(360\) sayısının kaç tane tek pozitif böleni vardır?
\item \(30!\) sayısının asal çarpanlarına ayrılışında \(5\)'in üssü kaçtır?
\item Bir doğal sayının \(6\) ile bölümünden kalan \(4\)'tür. Bu sayının \(5\) katının \(6\) ile bölümünden kalan kaçtır?
\item \(a\)'nın \(11\) ile bölümünden kalan \(6\), \(b\)'ninki \(9\)'dur. \(a+b\) ve \(a\cdot b\)'nin \(11\) ile bölümünden kalanlar kaçtır?
\item \(3^{50}\) sayısının birler basamağındaki rakam kaçtır?
\end{enumerate}

\subsection*{Cevap Anahtarı}
{\small
\begin{multicols}{2}
\begin{enumerate}[label=\textbf{\arabic*.},leftmargin=*,itemsep=1pt]
\item \(31,37,41,43,47,53,59\): \(7\) tane.
\item \(187=11\cdot 17\) asal değil; \(191\) asal.
\item \(2+29=31\Rightarrow a\cdot b=58\).
\item \(252=2^2\cdot 3^2\cdot 7\); \(18\) bölen.
\item \(400=2^4\cdot 5^2\); \(15\) pozitif, \(30\) tam bölen.
\item \((a+1)\cdot 4=20\Rightarrow a=4\).
\item \(1+2+3+6+9+18=39\).
\item \(360=2^3\cdot 3^2\cdot 5\); tek bölen \(3\cdot 2=6\).
\item \(\lfloor 30/5\rfloor+\lfloor 30/25\rfloor=6+1=7\).
\item \(5\cdot 4=20=18+2\Rightarrow\) kalan \(2\).
\item \(a+b\to 4\); \ \(a\cdot b\to 10\).
\item \(3,9,7,1\) tekrarı; \(50=4\cdot 12+2\Rightarrow 9\).
\end{enumerate}
\end{multicols}}
